\section{Generative Models}
What we have seen so far focused only on supervised learning tasks, where the goal was to learn a mapping from an input to an output. 
In this chapter, we will focus on unsupervised learning generative models, which are a class of machine learning models that aim to generate new data samples that are similar to the training data. 

In general, generative models can be divided into two main categories:
\begin{itemize}
    \item \textbf{Non-Probabilistic Models:} These models do not explicitly model the underlying probability distribution of the data. 
    Instead, they learn a mapping from a latent space to the data space. \textbf{Autoencoders} are a common example of non-probabilistic generative models.
    \item \textbf{Probabilistic Models:} These models explicitly model the underlying probability distribution of the data.
    The idea is to use a Neural Network to learn a mapping between a basic probability distribution (e.g., a Gaussian distribution) and a more complex distribution $p_{model}$.
    The model is trained to make the distribution as close as possible to the true data distribution $p_{data}$.
    \textbf{Generative Adversarial Networks (GANs)}, \textbf{Variational Autoencoders (VAEs)}, and \textbf{Diffusion Models} are examples of probabilistic generative models.
\end{itemize}

Generative models have a wide range of applications, including image generation, style transfer, and data augmentation.


Probabilistic Models can be further divided into two main categories:
\begin{itemize}
    \item \textbf{Explicit Density Models:} These models explicitly define a probability distribution over the data and learn the parameters of this distribution. 
    The main example is the \textbf{Variational Autoencoder (VAE)}.
    
    \item \textbf{Implicit Density Models:} These models do not explicitly define a probability distribution, but instead learn to generate samples from the data distribution.
    The main example is the \textbf{Generative Adversarial Network (GAN)}.
\end{itemize}

\subsection{GAN}
\textbf{Generative Adversarial Networks} (GANs) are actually a union of two neural networks rather than a single neural network.

In particular, GANs consist of:

\begin{itemize}
    \item \textbf{Generator:} The generator network takes a random noise vector as input and generates a sample that resembles the training data. 
    The goal of the generator is to produce samples that are indistinguishable from real data. 
    
    Because it is difficult to sample directly from a complex training distribution $p_{data}$, 
    the generator first samples from a simple distribution $p_{z}$ (such as a standard normal distribution). 
    It then learns to transform these simple samples into samples that mirror the training data distribution.
    
    \item \textbf{Discriminator:} The discriminator network takes a sample as input and outputs a probability $\in [0,1]$ 
    that the sample is real (from the training data) rather than fake (generated by the generator). 
    The discriminator acts both as a classifier and, implicitly, as a loss function for the generator network, with the goal of correctly classifying samples as real or fake.
\end{itemize}

\begin{comment}
    #TODO: add figure of GAN architecture
\end{comment}

\subsubsection{Adversarial Training and Objective Function}
The network is trained in an \textbf{adversarial manner}, where the generator and discriminator are trained simultaneously. 
If the discriminator gets better at distinguishing real from fake samples, the generator must improve to produce more realistic samples. 
This implies that GANs do not consider any explicit probability distribution, but instead focus on generating samples that are indistinguishable from real data.

Because we are dealing with this adversarial dynamic, the objective function is framed as a \textbf{min-max} game between the generator and discriminator networks. 
The objective function can be expressed as follows:
\begin{equation}
    \min_{\theta_G} \max_{\theta_D} \left[ \mathbb{E}_{x \sim p_{data}}[\underbrace{\log D_{\theta_D}(x)}_{\text{Discriminator output for real data}}] + \mathbb{E}_{z \sim p_{z}}[ \underbrace{\log(1 - D_{\theta_D}(G_{\theta_G}(z)))}_{\text{Discriminator output for fake data}}] \right]
\end{equation}

where $\theta_G$ and $\theta_D$ are the parameters of the generator and discriminator networks, respectively.

\begin{itemize}
    \item The \textbf{discriminator} maximizes the objective with respect to $\theta_D$ by assigning high probabilities to real samples, i.e., $D_{\theta_D}(x)\rightarrow 1$, and low probabilities to generated samples, i.e., $D_{\theta_D}(G_{\theta_G}(z))\rightarrow 0$.

    \item The \textbf{generator} minimizes the objective with respect to $\theta_G$ by generating samples $G_{\theta_G}(z)$ that maximize the discriminator's output, i.e., $D_{\theta_D}(G_{\theta_G}(z))\rightarrow 1$, making generated samples indistinguishable from real ones.
\end{itemize}

This adversarial approach has since been extended to many other applications, including image-to-image translation, text-to-image generation, and video generation.

Altough this a very strong theoretical framework, such objective function had 2 main problems in practice:
\begin{itemize}
    \item Such min-max game is very hard to optimize in practice, requiring an alternation of gradient ascents on the discriminator and gradient descents on the generator.
    
    \item The derivative of the objective function with respect to the generator's parameters heavily suffers from vanishing gradients. 
    In particular, when the discriminator is very good at distinguishing real from fake samples, the generator's gradients become very small, making it difficult for the generator to improve.
    Instead, gradient signals are very strong when the discriminator is poor at distinguishing real from fake samples, which can lead to unstable training dynamics.
\end{itemize}

To fix both problems, a common approach is to use a \textbf{non-saturating loss} for the generator, which maximizes $\log D_{\theta_D}(G_{\theta_G}(z))$ instead of minimizing $\log(1 - D_{\theta_D}(G_{\theta_G}(z)))$.
The resulting objective function is as follows:
\begin{equation}
    \max_{\theta_G} \max_{\theta_D} \left[ \mathbb{E}_{x \sim p_{data}}[\log D_{\theta_D}(x)] + \mathbb{E}_{z \sim p_{z}}[ \log(D_{\theta_D}(G_{\theta_G}(z)))] \right]
\end{equation}

\paragraph{Latent Space Arithmetic}
It was observed that the latent space of GANs exhibits interesting properties, such as the ability to perform arithmetic operations on the latent vectors.
For example, given a latent vector $z_1$ that generates an image of a smiling woman, a latent vector $z_2$ that generates an image of a neutral woman and a latent vector $z_3$ that generates an image of a neutral man, we can perform the following arithmetic operation in the latent space:
\begin{equation}
    z_4 = z_1 - z_2 + z_3
\end{equation}
The resulting latent vector $z_4$ can then be passed through the generator to produce an image of a smiling man. 
This property of the latent space allows for interesting manipulations of the generated samples, such as changing facial expressions or adding/removing objects from images.

\subsubsection{Inference}
After training, the generator can be used to generate new samples by sampling from the simple distribution $p_{z}$ and passing these samples through the generator network.
Usually, the discriminator is discarded after training, as it is no longer needed for generating new samples.

\subsubsection{Mode Collapse}
We would like the generator to learn to generate samples that cover the entire data distribution. Hower, in practice, earlier versions of GANs
suffered from a problem called \textbf{mode collapse}. The generator would learn to produce only a limited variety of samples, often collapsing to a single mode of the data distribution.

For example, if the training data consists of images of different animals, the generator might learn to produce only images of cats, ignoring other animals in the dataset.

\subsection{GAN Architectures}
What we have discussed so far is the basic GAN architecture, but we still need to discuss how these components are implemented in practice.
We will now discuss some of the most important GAN architectures that have been proposed in the literature.

\subsubsection{DCGAN}
\textbf{Deep Convolutional GANs} (DCGANs) are a class of GANs that use deep convolutional neural networks for both the generator and discriminator networks.

\paragraph{Generator:} Starting from a random noise vector, the generator network uses a series of convolutional layers to upsample the noise vector into a
final $64 \times 64 \times 3$ image.

\paragraph{Discriminator:} The discriminator network again uses a series of convolutional layers  to downsample the input image into a single scalar output.
To do so, it doesn't use any pooling layers, but instead uses strided convolutions to downsample the input image.

\subsubsection{LSGAN}
\textbf{Least Squares GANs} (LSGANs) are a class of GANs that replace the binary cross-entropy loss used in the original GAN formulation with a least squares loss.
This change had 2 main advantages:
\begin{itemize}
    \item Least squares loss provides stronger gradients for the generator when the discriminator is confident, which helps to mitigate the vanishing gradient problem
    we discussed earlier.
    \item The stronger gradients also allow the generator to escape from local minima, mitigating the mode collapse problem we discussed earlier.
\end{itemize}

After this work, many other GAN variants have been proposed, each with its own advantages and disadvantages.

\subsubsection{Wasserstein GAN (WGAN)}
\textbf{Wasserstein GANs} (WGANs) try to further improve the training stability of LSGANs by using the Wasserstein distance. 
The original GAN formulation is equivalent to minimizing the Jensen-Shannon divergence between the real and generated data distributions.
However, the Jensen-Shannon divergence has been shown to be a poor choice for training GANs, as it can lead to vanishing gradients and mode collapse.

The Wasserstein distance, on the other hand, is a more robust metric for measuring the distance between two probability distributions.
The Wasserstein distance has an important property: it changes smoothly, even when the two distributions do not overlap.
This property allows the generator to receive meaningful gradients even when the discriminator is very confident, which helps to mitigate the vanishing gradient problem and improve training stability.

\subsection{Conditioning the Generation}
All the architectures we have discussed so far are unconditional, meaning that the generator produces samples without any control over the output.
GANS can be conditioned on additional information, such as class labels, text descriptions, or even other images, to generate samples that satisfy specific conditions.

\subsubsection{Conditional GAN}
\textbf{Conditional GANs} (cGANs) are a class of GANs that allow for conditioning the generation process trough labels.
In particular, the generator takes as input both a random noise vector and a label, and generates a sample that corresponds to the given label.
The discriminator also takes as input a sample and a label for both real and generated samples, and outputs a probability that the sample is real given the label.

The objective function for cGANs is similar to the original GAN formulation, but with the addition of the conditioning label:
\begin{equation}
    \min_G \max_D V(D, G) = \mathbb{E}_{x \sim p_{data}}[\log D(x | y)] + \mathbb{E}_{z \sim p_z}[\log(1 - D(G(z) | y))]
\end{equation}

\subsubsection{Text-To-Image}
Following the success of cGANs, several works have proposed to condition the generation process on text descriptions, allowing for \textbf{text-to-image synthesis}.
In this setting, some text description is appended to the random noise vector and passed through the generator, which produces an image that corresponds to the given text description.

Similarly to the cGANs, the discriminator also takes as input both an image and a text description, and outputs a probability that the image is real given the text description.

\subsubsection{Image-to-Image Translation}
One of the most popular applications of GANs is \textbf{image-to-image translation}, where the goal is to learn a mapping from one image domain to another.
For example, given a dataset of black-and-white images and their corresponding color images, a GAN can be trained to learn the mapping from black-and-white to color images.

Such architecture is called \textbf{Pix2Pix}, and it uses a CNN-based generator that instead of generating an image from a random noise vector, takes an input image and generates a corresponding output image,
of the same size.
Again, the discriminator takes as input both the input image and the generated output image, and outputs a probability that the output image is real given the input image.

\subsubsection{CycleGAN}
\textbf{CycleGAN} is a GAN architecture that allows for image-to-image translation without paired training data. The previously discussed Pix2Pix architecture requires paired training data,
meaning that for each input image, there must be a corresponding output image in the training dataset.
However, in many cases, paired training data is not available.

The architecture of CycleGAN consists of two generators and two discriminators. 
The first generator $F$ learns to map images from the source domain $X$ to the target domain $Y$, generating $\hat{y}$ from $x$.
The second generator $G$ learns to map images from the target domain $Y$ back to the source domain $X$, generating $\hat{x}$ from $\hat{y}$.
The discriminators take as input both the generated image and a real image from their respective domains, and output probabilities that the generated images are real given the input images.

To ensure that the generators learn meaningful mappings, CycleGAN introduces a \textbf{cycle consistency loss}, which encourages the generators to produce images that can be mapped back to the original domain.
This can be expressed as follows:
\begin{equation}
    \mathcal{L}_{cycle}(F, G) = \mathbb{E}_{x \sim p_{data}(x)}[\|G(F(x)) - x\|_1] + \mathbb{E}_{y \sim p_{data}(y)}[\|F(G(y)) - y\|_1]
\end{equation}

\subsubsection{StyleGAN}
\textbf{StyleGAN} is a GAN architecture proposed by NVIDIA that allows for fine-grained control over the generated images by introducing a style-based generator architecture.
To enable this, StyleGAN introduces a \textbf{mapping network} that thakes the random noise vector as input and maps it to an intermediate latent space, with better features for controlling the style of the generated images.

The intermediate latent vector $w$ is then passed through a series of \textbf{adaptive instance normalization} (AdaIN) layers, which control the style of the generated image at different levels of detail. 
This allows for controlling the overall structure of the image, as well as finer details such as textures and colors.

\paragraph{AdaIN}
Adaptive Instance Normalization (AdaIN) is a normalization technique taken from style transfer literature. 
Technically, this is an instance normalization layer that instead of using learned parameters from the training data, 
uses the parameters computed from the style image, which in the case of StyleGAN is the intermediate latent vector $w$.

\subsection{Training Tricks}
A few training tricks have been proposed to improve the training stability of GANs.

\begin{itemize}
    \item \textbf{Minibatch Discrimination:} used to reduce mode collapse. It allows the discriminator to look at
    the similarity of an image $x$ to other images in the same minibatch. The similarity is passed to the discriminator as an additional feature,
    allowing it to detect if the generator is producing a limited variety of samples.
    \item \textbf{Virtual Batch Normalization:} Due to batch normalization, the generated images are not independent of each other.
    To mitigate this problem, virtual batch normalization uses a reference batch of real images to compute the normalization statistics for each generated image,
    instead of using the statistics of the current minibatch.
\end{itemize}

\subsection{Evaluating GANs Performance}
Since the inception of GANs, evaluating their performance has been a challenging task. To this day, the evaluation of generative models remains an open research problem, and there is no single universally accepted metric to assess GAN performance.

Early approaches relied on human evaluators to assess the quality of generated samples. However, this method is unscalable, expensive, and highly subjective, as different annotators may have conflicting opinions on visual quality.

Additionally, computing the likelihood of generated samples or directly comparing the real and generated data distributions is intractable, as GANs do not explicitly model the underlying probability distribution of the data.

\subsubsection{Inception Score (IS)}
The first metric proposed to replace human evaluation was the \textbf{Inception Score} (IS). The core idea is that a good generator should produce highly recognizable objects from a wide variety of classes (meaning the generated sample classes should be uniformly distributed).

To compute the IS, a pre-trained Inception model is used to classify the generated samples. The score evaluates two criteria:
\begin{enumerate}
    \item The generated images should be easily classifiable, meaning the predicted class distribution $p(y|x)$ has \textbf{low entropy}.
    \item The generator should output a diverse range of images, meaning the marginal class distribution over all generated samples $p(y)$ has \textbf{high entropy}.
\end{enumerate}

Formally, the Inception Score is computed as follows:
\begin{equation}
    IS = \exp(\mathbb{E}_{x \sim p_{g}}[D_{KL}(p(y|x) || p(y))])
\end{equation}
This formula calculates the \textbf{Kullback-Leibler divergence} between the conditional label distribution $p(y|x)$ and the marginal label distribution $p(y)$. A higher KL divergence results in a higher (better) Inception Score, which occurs when the network is highly confident about individual images but predicts a uniform distribution of classes overall.

\paragraph{Issues with the Inception Score}
The Inception Score suffers from two main issues:
\begin{itemize}
    \item \textbf{Memorization:} A GAN that simply memorizes the training data will achieve a high IS, as the generated samples will perfectly match real images (easily classifiable and evenly distributed).
    \item \textbf{Mode Collapse:} If the generator produces only a single realistic sample for each class, the IS can still be high. The samples are evenly distributed across classes and easily recognized, but the model entirely lacks intra-class diversity.
\end{itemize}

\subsubsection{Frechet Inception Distance (FID)}

The \textbf{Fréchet Inception Distance} (FID) was proposed to overcome several shortcomings of the Inception Score. 
Rather than relying solely on the predicted class probabilities, FID compares the distributions of feature representations extracted from real and generated samples.

To compute FID, both real and generated samples are passed through a pre-trained Inception network, and feature vectors are extracted from one of its intermediate layers. 
The feature distributions are then approximated as multivariate Gaussian distributions, characterized by their means and covariance matrices.

The FID score is defined as
\begin{equation}
    \mathrm{FID} =
    \|\mu_r-\mu_g\|_2^2 +
    \mathrm{Tr}\!\left(
        \Sigma_r+\Sigma_g
        -2(\Sigma_r\Sigma_g)^{1/2}
    \right),
\end{equation}
where $(\mu_r,\Sigma_r)$ and $(\mu_g,\Sigma_g)$ denote the mean vectors and covariance matrices of the real and generated feature distributions, respectively.

This expression corresponds to the Fréchet distance between the two Gaussian distributions.
A lower FID score indicates that the generated feature distribution is closer to that of the real data, suggesting higher visual quality and greater diversity of the generated samples.

Compared to IS, FID directly measures the similarity between generated and real data distributions. 
Nevertheless, it still has limitations. In particular, FID cannot explicitly detect overfitting, as a generator that simply memorizes the training set would achieve an FID score close to zero.

\subsubsection{Kernel Inception Distance (KID)}

Like FID, the \textbf{Kernel Inception Distance} (KID) evaluates the similarity between the feature distributions of real and generated samples extracted from a pre-trained Inception network.

Instead of computing the Fréchet distance between Gaussian approximations of the feature distributions, KID measures their discrepancy using the \textbf{Maximum Mean Discrepancy} (MMD). 
Specifically, KID computes the squared MMD between the two sets of feature vectors using a polynomial kernel.

A lower KID score indicates that the generated feature distribution is more similar to the real one.

Compared to FID, the main advantage of KID is that its estimator is \textbf{unbiased}, whereas the empirical estimator of FID is biased, particularly when computed from a limited number of samples. 
Consequently, KID generally provides a more reliable estimate of distribution similarity for smaller datasets, while maintaining comparable performance on larger datasets.

\subsubsection{Precision and Recall}

The well-known \textbf{precision} and \textbf{recall} metrics have also been extended to evaluate generative models.

In this setting, \textbf{precision} measures the quality (or fidelity) of the generated samples by estimating the fraction of generated samples that lie within the support of the real data distribution. 
Conversely, \textbf{recall} measures the diversity (or coverage) of the generated distribution by estimating the fraction of real samples that lie within the support of the generated distribution.

Pratically, precision and recall are computed by embedding both real and generated samples into a feature space using a pre-trained network (e.g., Inception).
The support of each distribution is then estimated using nearest neighbor methods.

\subsubsection{Perceptual Path Length}

The \textbf{Perceptual Path Length} (PPL) metric was introduced to evaluate the quality of the latent space learned by a generative model.
In particular, it measures the smoothness of the latent space by quantifying how small changes in the latent vector affect the generated samples.

Intuitively, picking two points in the latent space and interpolating between them should result in a smooth transition between the corresponding generated samples.

\subsubsection{Issues with GAN Evaluation Metrics}
Despite the variety of metrics proposed to evaluate GANs, there is still no universally accepted standard that comprehensively assesses the performance of generative models. Each metric possesses distinct strengths and limitations, often capturing entirely different aspects of the generated distribution. 

A fundamental challenge in evaluating GANs is that the generator's loss does not directly correlate with the visual or structural quality of the generated samples. In traditional Machine Learning, a core principle is that the loss function serves as a reliable proxy for model performance. However, in the adversarial framework of a GAN, the generator's loss is a moving target. Because it is heavily influenced by—and entirely relative to—the discriminator's current performance, a low generator loss simply indicates that it is currently fooling the discriminator, not necessarily that it is producing high-quality, realistic samples.

